\section*{Capítulo \ref{ch:b3:petrological-thermodynamics} --- Termodinâmica das rochas}

\begin{solution}{exo:b3:petrological-thermodynamics:1}
$\Delta_rH = -3029.3 - 910.7 + 3935.0 = \qty{-5.0}{kJ}$; $\Delta_rS = 133.5 +
41.5 - 207.4 = \qty{-32.4}{J.K^{-1}}$; $\Delta_rV = 60.40 + 22.69 - 100.07 =
\qty{-16.98}{cm^3} = \qty{-1.698}{J.bar^{-1}}$. Inclinação $32.4/1.698 =
\qty{19.1}{bar/K}$. A \qty{873}{K}: $P = 1 + (873.15\times(-32.4) +
5000)/(-1.698) = \qty{13.7}{kbar}$.
\end{solution}

\begin{solution}{exo:b3:petrological-thermodynamics:2}
$c = 2$ e $p = 3 > c$: a regra é violada. As três fases coexistem somente sobre a
linha univariante de \ce{Mg2SiO4 + SiO2 <=> 2 MgSiO3}; uma rocha coletada ao
acaso não está ali, de modo que a paragênese registra uma reação incompleta,
por exemplo forsterita envolta por uma borda de enstatita que a impediu de
reagir com o quartzo.
\end{solution}

\begin{solution}{exo:b3:petrological-thermodynamics:3}
$a_{\text{grossular}} = 0.20^3 = 0.008$; $a_{\text{piropo}} = 0.15^3 =
\num{3.4e-3}$; $a_{\text{anortita}} = 0.40$.
\end{solution}

\begin{solution}{exo:b3:petrological-thermodynamics:4}
$T_c = 15\,000/(2\times8.314) = \qty{902}{K}$; $W = 2\times8.314\times1000
= \qty{16.6}{kJ.mol^{-1}}$.
\end{solution}

\begin{solution}{exo:b3:petrological-thermodynamics:5}
$\Delta_rG = -5432.2 + 2\times(-2088.8) - 730.2 + 3\times3478.1 =
\qty{94.3}{kJ}$; $P = 1 + 94\,300/6.61 = \qty{14.3}{kbar}$. A linha
experimental dá $22.80\times726.85 - 1093 = \qty{15.5}{kbar}$, \qty{1.2}{kbar}
acima.
\end{solution}

\begin{solution}{exo:b3:petrological-thermodynamics:6}
$RT = \qty{7483}{J.mol^{-1}}$. Para $x = 0.5$: $7483\ln0.5 + 20\,000/4 = -5187 +
5000 = \qty{-187}{J.mol^{-1}}$. Para $x = 0.111$: $7483\times(0.111\ln0.111 +
0.889\ln0.889) + 20\,000\times0.111\times0.889 = -2609 + 1974 =
\qty{-635}{J.mol^{-1}}$, e o mesmo valor em 0.889 por simetria. Quantidades iguais
dos dois flancos têm composição global 0.5 e energia de Gibbs
\qty{-635}{J.mol^{-1}}: a separação a reduz em \qty{448}{J.mol^{-1}}.
\end{solution}

\begin{solution}{exo:b3:petrological-thermodynamics:7}
$c + 2 = 5$ fases e 5 linhas. (D) $A + C = 2B$: $0 + 2/3 = 2\times1/3$. (A)
$B + D = 2C$: $1/3 + 1 = 2\times2/3$. (C) $2A + D = 3B$: $0 + 1 = 3\times1/3$.
(B) $A + 2D = 3C$: $0 + 2 = 3\times2/3$.
\end{solution}

\begin{solution}{exo:b3:petrological-thermodynamics:8}
Pela regra da alavanca, a fase em 0.889 corresponde a $(0.3 - 0.111)/(0.889 - 0.111) =
0.243$, \qty{24}{\percent}, e a fase em 0.111, \qty{76}{\percent}.
\end{solution}

\begin{solution}{exo:b3:petrological-thermodynamics:9}
$\partial P/\partial X_{\text{An}} = -3RT/(6.21X_{\text{An}}) =
-3\times8.314\times900/(6.21\times0.30) = \qty{-12000}{bar}$ por unidade de
$X_{\text{An}}$: um erro de 0.01 desloca a pressão em \qty{120}{bar}.
\end{solution}

\begin{solution}{exo:b3:petrological-thermodynamics:10}
$G'_{\text{mix}} = RT\ln[x/(1 - x)] + W(1 - 2x)$ e $G''_{\text{mix}} =
RT/[x(1 - x)] - 2W$, nulo onde $x(1 - x) = RT/(2W)$. Para $x = 0.3$: $T =
2\times20\,000\times0.21/8.314 = \qty{1010}{K}$; o \omterm{def:b3:petrological-thermodynamics:solvus}{solvus} é atingido em
$20\,000\times0.4/(8.314\ln(7/3)) = \qty{1136}{K}$. Entre os dois valores, o
cristal é metaestável: a separação reduziria sua energia, mas uma pequena
flutuação a eleva, de modo que a segunda fase precisa nuclear.
\end{solution}

\begin{solution}{exo:b3:petrological-thermodynamics:11}
Como $G_{\text{mix}}(x) = G_{\text{mix}}(1 - x)$, a tangente comum toca a curva
em dois pontos simétricos $x$ e $1 - x$ e é horizontal:
$G'_{\text{mix}}(x) = 0$, que é a equação. Resolvendo, $T(x) = W(1 -
2x)/[R\ln((1 - x)/x)]$. Com $x = 1/2 - \varepsilon$, $1 - 2x = 2\varepsilon$
e $\ln((1/2 + \varepsilon)/(1/2 - \varepsilon)) \approx 4\varepsilon$, logo
$T \to W/(2R)$ quando $\varepsilon \to 0$; e, para $x < 1/2$, $T(x)$ cresce à
medida que $x$ se aproxima de 1/2, de modo que esse limite é o máximo.
\end{solution}

\begin{solution}{exo:b3:petrological-thermodynamics:12}
$(2.0 - 0.08\times10)/0.92 = \qty{1.30}{wt\percent}$. O manganês favorece a
granada: com menos dele na \omterm{def:b3:petrological-thermodynamics:pseudosection}{composição global efetiva}, as bordas crescem em um
campo diferente e ao longo de isopletas distintas daquelas que a análise da
rocha total preveria.
\end{solution}

\begin{solution}{pb:b3:petrological-thermodynamics:1}
\textbf{1.} $\Delta_rH^\circ = -6640.0 + 2\times(-2593.8) - 910.7 + 3\times4234.0
= \qty{-36.3}{kJ}$; $\Delta_rS^\circ = 260.1 + 165.6 + 41.5 - 597.9 =
\qty{-130.7}{J.K^{-1}}$; $\Delta_rV = 125.28 + 88.30 + 22.69 - 302.37 =
\qty{-66.10}{cm^3} = \qty{-6.61}{J.bar^{-1}}$.

\textbf{2.} $-36\,300 + 298.15\times130.7 = \qty{+2.67}{kJ}$: positivo, logo
a anortita é estável na superfície.

\textbf{3.} $P = 1 + (900\times(-130.7) + 36\,300)/(-6.61) = \qty{12.3}{kbar}$;
inclinação $130.7/6.61 = \qty{19.8}{bar/K}$.

\textbf{4.} $1 + 81\,500/6.61 = \qty{12.3}{kbar}$ a \qty{900}{K}; $1 +
119\,800/6.61 = \qty{18.1}{kbar}$ a \qty{1200}{K}.

\textbf{5.} A 626.85 e \qty{926.85}{\celsius} a linha experimental dá
13.2 e \qty{20.0}{kbar}: 0.9 e \qty{1.9}{kbar} acima das tabelas.

\textbf{6.} $1914\times6.61 = \qty{12.7}{kJ}$. A incerteza combinada é
$\sqrt{3.2^2 + (3\times2.0)^2 + (2\times2.0)^2 + 1.0^2} = \qty{8.0}{kJ}$: a
diferença é 1.6 vez esse valor, compatível com calorimetria, e é a razão de a
linha ser tirada dos experimentos.

\textbf{7.} $\Delta_rC_p$ é pequeno entre os sólidos, e suas duas integrais se
cancelam em parte: as tabelas deslocam a linha em apenas 26 e \qty{112}{bar}.

\textbf{8.} \ce{SiO2}, \ce{Al2O3}, \ce{FeO}, \ce{MgO}, \ce{MnO}, \ce{CaO},
\ce{Na2O}, \ce{K2O} e \ce{H2O}: $c = 9$. Fases: granada, biotita,
 muscovita, plagioclásio, cianita, quartzo e fluido, $p = 7$.

\textbf{9.} $F = 9 - 7 + 2 = 4$: pressão, temperatura e duas
variáveis composicionais. $p = 7 \le c = 9$: a regra vale, e a associação pode ser estável
em uma área do plano pressão--temperatura.

\textbf{10.} Dois polimorfos puros de \ce{Al2SiO5} em equilíbrio situam-se sobre a
linha univariante cianita = sillimanita: uma relação entre $P$ e $T$, e um
termômetro fornece então a pressão.

\textbf{11.} O fluido impõe o potencial químico da água: 8 componentes livres
para 6 fases sólidas, $F = 8 - 6 + 2 = 4$, e $p \le c$ ainda
vale.

\textbf{12.} $a_{\text{grossulária}} = 0.10^3 = 0.001$, $a_{\text{anortita}} =
0.30$; $\ln K = \ln(0.001/0.30^3) = -3.30$.

\textbf{13.} $K_D = (0.12/0.74)/(0.45/0.55) = 0.1622/0.8182 = 0.198$.

\textbf{14.} $T = (52\,635 + 2488)/(20.20 + 24.94\times1.618) = 55\,123/60.57
= \qty{910}{K}$, ou \qty{637}{\celsius}.

\textbf{15.} $P_0 = 22.80\times636.9 - 1093 = \qty{13429}{bar}$; $(3\times8.314
\times910.1/6.21)\ln(1/3) = \qty{-4016}{bar}$; $P = \qty{9.41}{kbar}$.

\textbf{16.} Em \qty{9.41}{kbar}, \qty{644.2}{\celsius}, em seguida
\qty{9.55}{kbar}; nessa pressão, \qty{644.9}{\celsius} e
\qty{9.56}{kbar}, sem alteração depois: \qty{645}{\celsius} e
\qty{9.56}{kbar}.

\textbf{17.} \qty{641}{\celsius} e \qty{8.79}{kbar}. As tabelas colocam a
linha do membro final cerca de \qty{1}{kbar} baixa demais (Parte I); a linha experimental
é mantida.

\textbf{18.} Com $X_{\text{an}} = 0.35$: \qty{642}{\celsius} e
\qty{8.93}{kbar}, \qty{0.63}{kbar} a menos.

\textbf{19.} A \qty{918}{K}: $P = 1 + (918\times12.6 - 7700)/0.571 =
\qty{6.8}{kbar}$, abaixo de \qty{9.56}{kbar}: cianita, como observado.

\textbf{20.} $z = \num{9.56e8}/(2800\times9.81) = \qty{34.8}{km}$.

\textbf{21.} $\Delta P = (3\times8.314\times918/6.21)\ln(0.07/0.10) =
\qty{-1.3}{kbar}$: o núcleo se formou cerca de \qty{4.8}{km} mais raso. A granada
cresceu durante o soterramento, e sua borda registra o ponto mais profundo.

\textbf{22.} O xisto Varden atingiu uma pressão máxima de
\textbf{\qty[text-series-to-math=true]{9.6}{kbar}} a
\qty{645}{\celsius}, com um soterramento de cerca de \qty{35}{km}.
\end{solution}
